\documentclass[11pt,a4paper]{article}

\usepackage{fontspec}
\usepackage{polyglossia}
\setdefaultlanguage{russian}
\setotherlanguage{english}
\setmainfont{DejaVu Serif}
\setsansfont{DejaVu Sans}
\setmonofont{DejaVu Sans Mono}

\usepackage[a4paper,margin=24mm,headheight=15pt]{geometry}
\usepackage{amsmath,amssymb,amsthm,mathtools}
\usepackage{microtype}
\usepackage{booktabs,longtable,tabularx,array}
\usepackage{enumitem}
\usepackage{xcolor}
\usepackage{hyperref}
\usepackage{fancyhdr}
\usepackage{listings}
\usepackage{fvextra}
\usepackage{tikz}
\usetikzlibrary{arrows.meta,positioning,calc,fit}

\hypersetup{
  colorlinks=true,
  linkcolor=blue!55!black,
  urlcolor=blue!55!black,
  citecolor=blue!55!black,
  pdftitle={Терминальная PBW-импликация и план завершения Lean-формализации},
  pdfauthor={}
}

\pagestyle{fancy}
\fancyhf{}
\lhead{Терминальная PBW-импликация}
\rhead{План Lean-формализации}
\cfoot{\thepage}

\setlist[itemize]{leftmargin=2em,itemsep=0.35em,topsep=0.45em}
\setlist[enumerate]{leftmargin=2.2em,itemsep=0.4em,topsep=0.45em}
\setlength{\parindent}{0pt}
\setlength{\parskip}{0.55em}
\allowdisplaybreaks

\newtheorem{theorem}{Теорема}[section]
\newtheorem{lemma}[theorem]{Лемма}
\newtheorem{proposition}[theorem]{Предложение}
\newtheorem{corollary}[theorem]{Следствие}
\theoremstyle{definition}
\newtheorem{definition}[theorem]{Определение}
\newtheorem{construction}[theorem]{Конструкция}
\theoremstyle{remark}
\newtheorem{remark}[theorem]{Замечание}
\newtheorem{warning}[theorem]{Предупреждение}

\newcommand{\Z}{\mathbb Z}
\newcommand{\Q}{\mathbb Q}
\newcommand{\Sym}{\operatorname{Sym}}
\newcommand{\pr}{\operatorname{pr}}
\newcommand{\Rel}{\mathcal R}
\newcommand{\Obs}{\operatorname{Obs}}
\newcommand{\im}{\operatorname{im}}
\newcommand{\coker}{\operatorname{coker}}
\newcommand{\ev}{\operatorname{ev}}
\newcommand{\front}{\operatorname{Front}}
\newcommand{\src}{\operatorname{src}}
\newcommand{\out}{\operatorname{out}}
\newcommand{\wt}{\operatorname{wt}}
\newcommand{\PBW}{\operatorname{PBW}}
\newcommand{\id}{\operatorname{id}}
\newcommand{\barF}{\overline F}
\newcommand{\eps}{\varepsilon}

\definecolor{darkred}{RGB}{130,20,20}
\definecolor{darkgreen}{RGB}{20,100,45}
\definecolor{lightgray}{RGB}{245,245,245}
\definecolor{lightblue}{RGB}{238,246,255}
\definecolor{lightred}{RGB}{255,242,242}
\definecolor{lightgreen}{RGB}{240,250,242}

\newenvironment{keybox}{%
  \par\medskip\noindent
  \begin{center}\begin{minipage}{0.94\textwidth}
  \color{blue!55!black}\hrule\vspace{0.55em}
}{%
  \vspace{0.35em}\hrule\end{minipage}\end{center}\medskip
}

\newenvironment{dangerbox}{%
  \par\medskip\noindent
  \begin{center}\begin{minipage}{0.94\textwidth}
  \color{darkred}\hrule\vspace{0.55em}
}{%
  \vspace{0.35em}\hrule\end{minipage}\end{center}\medskip
}

\newenvironment{successbox}{%
  \par\medskip\noindent
  \begin{center}\begin{minipage}{0.94\textwidth}
  \color{darkgreen}\hrule\vspace{0.55em}
}{%
  \vspace{0.35em}\hrule\end{minipage}\end{center}\medskip
}
\lstdefinelanguage{Lean}{
  morekeywords={def,theorem,lemma,structure,where,extends,inductive,match,with,if,then,else,let,in,by,have,show,exact,apply,refine,constructor,termination_by,decreasing_by,namespace,end,import,open,private,protected,instance,deriving,forall,exists},
  sensitive=true,
  morecomment=[l]{--},
  morecomment=[s]{/-}{-/},
  morestring=[b]"
}
\lstset{
  language=Lean,
  basicstyle=\ttfamily\small,
  keywordstyle=\bfseries,
  commentstyle=\itshape,
  stringstyle=\ttfamily,
  backgroundcolor=\color{lightgray},
  frame=single,
  rulecolor=\color{black!20},
  breaklines=true,
  columns=fullflexible,
  keepspaces=true,
  showstringspaces=false,
  xleftmargin=0.4em,
  xrightmargin=0.4em,
  aboveskip=0.7em,
  belowskip=0.7em
}

\title{\textbf{Терминальная PBW-импликация}\\[0.35em]
\large Полная размеченная трасса, balanced ledger, нулевой obstruction\\
и план завершения формализации в Lean}
\author{}
\date{}

\begin{document}
\maketitle

\begin{abstract}
Настоящий документ изолирует и подробно доказывает сложную импликацию,
используемую в терминальной части PBW-сборки исходного доказательства.
Главный результат имеет вид
\[
 \text{полная размеченная PBW/truncation-трасса}
 \Longrightarrow
 \text{balanced ledger}
 \Longrightarrow
 \Obs(\mathcal W_{\rm exc})=0.
\]
Именно эта импликация была потеряна при попытке формализации: агент оставил
только абстрактную signed family $\mathcal W_{\rm exc}$ и пытался вывести
нулевой obstruction из принадлежности порождённому отношениями идеалу. Такое
утверждение неверно в общем случае. В исходном доказательстве имеется
дополнительный свидетель --- полная размеченная трасса всех элементарных
перестановок и усечений. Она сохраняет происхождение каждого occurrence,
все lower-context descendants и общий набор коэффициентов. Из неё строится
один и тот же элемент $z$, одновременно удовлетворяющий квадратичному и
примитивному равенствам.

После математического доказательства дан конкретный план Lean-реализации:
какую сессию использовать, какие данные нельзя стирать, какие определения и
леммы создать, как построить proof-producing collector и по каким критериям
считать задачу завершённой.
\end{abstract}

\tableofcontents
\newpage

\section{Итоговый диагноз}

\begin{keybox}
\textbf{Математический итог.}
Проблемная импликация в исходном доказательстве доказана. Её доказательство
не является следствием одного лишь факта
$\Theta\in\Rel U(F)$ и не является wordwise-утверждением. Оно использует
полную размеченную трассу преобразований. Если трассу сохранить как явный
свидетель, то из неё конструктивно получается balanced ledger, один общий
элемент $z$ и, следовательно, нулевой terminal obstruction.
\end{keybox}

\begin{dangerbox}
\textbf{Ошибка формализационной архитектуры.}
Агент пытался доказать утверждение примерно следующего вида:
\[
  \text{SignedRelationWordFamily}(\mathcal W)
  \;\Longrightarrow\;
  \Obs(\mathcal W)=0,
\]
используя в основном значение суммы $\Theta_{\mathcal W}\in\Rel U(F)$.
Это слишком сильное и вообще неверное утверждение. Правильная цель имеет
дополнительный аргумент:
\[
  \text{CompleteMarkedTrace}(\mathcal W)
  \;\Longrightarrow\;
  \Obs(\mathcal W)=0.
\]
Трасса не является технической декорацией: она и есть недостающий witness.
\end{dangerbox}

Формально ситуация состоит из трёх уровней.

\begin{enumerate}
\item \textbf{Локальный уровень.} Каждая PBW-перестановка, перестановка
обычного фактора с полным отношением и каждый шаг усечения являются буквальным
целочисленным тождеством.

\item \textbf{Глобальный уровень.} Все локальные шаги собираются в конечную
размеченную трассу. Prefix-ledger invariant доказывает точное равенство между
начальной signed family и конечным frontier.

\item \textbf{Терминальный уровень.} После полного, а не преждевременного,
суммирования frontier раскладывается на терминальные строки, целые relation
leaves и попарно согласованные горизонтальные/вертикальные диагонали. Из
терминальных строк определяется один элемент $z$; обе требуемые координаты
читаются с тех же самых occurrences и тех же коэффициентов.
\end{enumerate}

Именно переход от первого уровня сразу к терминальному, с забыванием второго,
создал ложную недостающую задачу.

\section{Положение импликации в исходном доказательстве}

Пусть $F$ --- градуированное относительно свободное метабелево кольцо Ли
класса не выше $n+1$,
\[
 F=\bigoplus_{s=1}^{n+1}F_s,
 \qquad M=F_{\ge 2}=[F,F],
 \qquad [M,M]=0,
\]
$\Rel\triangleleft F$ --- идеал отношений, а $U(F)$ --- универсальная
обёртывающая алгебра. Пусть
\[
 \Theta\in \Rel U(F)
\]
удовлетворяет governing conditions:

\begin{enumerate}[label=(G\arabic*)]
\item все многофакторные PBW-коэффициенты $\Theta$ веса не выше $2n$ равны
нулю;
\item однофакторная часть $\Theta$ однородна веса $n+1$.
\end{enumerate}

Для $2\le k\le n$ положим
\[
 A_k=F_{\le k},\qquad D_k=\pr_{\le k}(\Rel),
 \qquad m_k=n+2-k,
\]
и рассмотрим Koszul-комплексы
\[
 K_{m_k}(k)=K_{m_k}(D_k\hookrightarrow A_k).
\]
Точная PBW-сборка строит цепи
\[
 \chi_k\in K_{m_k}(k)_1\quad(2\le k<n),
 \qquad
 \chi_n\in Z_1K_2(n).
\]
Их границы являются полными factor-$m_k$ symbols элемента $\Theta$ после
проекции в $\Sym^{m_k}A_k$.

Обозначим через $Q_n(\chi_n)\in C$ ориентированный примитивный коэффициент
квадратичного блока и определим
\begin{equation}\label{eq:eps-def}
 \eps=
 \iota_C\bigl(\pr_{1,n+1}\Theta\bigr)
 -\sum_{k=2}^{n-1}\Phi_k(f_k\chi_k)
 -Q_n(\chi_n).
\end{equation}
Проблемный переход состоит в доказательстве равенства
\begin{equation}\label{eq:terminal-discrepancy}
 \eps=-\sum_{k=2}^{n-1}T_k(\partial\chi_k).
\end{equation}
Затем governing conditions дают
\begin{equation}\label{eq:T-vanish}
 T_k(\partial\chi_k)=0,
\end{equation}
и потому $\eps=0$.

Эквивалентная standalone-формулировка использует terminal obstruction. Пусть
$\pi:F\twoheadrightarrow A$ --- нужная проекция, $D=\pi(\Rel)$,
\[
 K=D\otimes A,
 \qquad B=\Sym^2A,
 \qquad \partial(d\otimes a)=d\odot a,
\]
а $\lambda:D\to\Rel$ и $s:A\to F$ --- аддитивные сечения. Определим
\[
 \sigma(d\otimes a)
 =P_1\bigl(\iota(\lambda d)\iota(sa)\bigr).
\]
Для signed full-relation word family $\mathcal W$ положим
\[
 \Theta_{\mathcal W}=\sum_e m_e\iota(\rho_e)X_e,
 \quad
 b_{\mathcal W}=\Sym^2(\pi)P_2(\Theta_{\mathcal W}),
 \quad
 p_{\mathcal W}=P_1(\Theta_{\mathcal W}).
\]
Далее
\[
 \Lambda:K\oplus\Rel\longrightarrow B\oplus F,
 \qquad
 \Lambda(z,r)=(\partial z,\sigma z-r),
\]
и
\[
 \Obs(\mathcal W)
 =[(b_{\mathcal W},p_{\mathcal W})]\in\coker\Lambda.
\]
По определению
\begin{equation}\label{eq:obstruction-criterion}
 \Obs(\mathcal W)=0
 \quad\Longleftrightarrow\quad
 \exists z\in D\otimes A\;\exists r\in\Rel:\;
 \begin{cases}
  \partial z=b_{\mathcal W},\\
  \sigma z=p_{\mathcal W}+r.
 \end{cases}
\end{equation}
Критически важно, что в обеих строках стоит \emph{один и тот же} $z$.

\section{Точный тезис, который требуется формализовать}

Обозначим через $E=\mathcal W_{\rm exc}$ exceptional upper-edge family,
возникающую внутри полной PBW/truncation-трассы. Правильный тезис имеет
следующий вид.

\begin{theorem}[Trace-to-obstruction implication]\label{thm:trace-to-obs}
Пусть дана полная конечная размеченная трасса элементарных PBW- и
truncation-преобразований, корневая signed family которой равна $E$. Пусть
трасса:
\begin{enumerate}[label=(T\arabic*)]
\item сохраняет каждое полное отношение как единый relation occurrence;
\item включает все descendants каждого корневого occurrence;
\item использует только буквальные transfer- и truncation-тождества;
\item заканчивается frontier, для которого выполнена исчерпывающая
классификация по числу факторов и весу;
\item сохраняет отдельные labels даже для occurrences с одинаковым
алгебраическим значением.
\end{enumerate}
Тогда $E$ допускает balanced ledger. В частности,
\[
 \Obs(E)=0.
\]
Более точно, из трассы конструктивно определяется один элемент
$z\in D\otimes A$ и элемент $r\in\Rel$, для которых
\[
 \partial z=b_E,
 \qquad
 \sigma z=p_E+r.
\]
\end{theorem}

Вся сложность состоит в первой импликации
\[
 \text{CompleteMarkedTrace}(E)
 \Longrightarrow
 \text{BalancedLedger}(E).
\]
Вторая импликация является немедленным применением двух additive reads к
ledger identity.

\section{Почему данные одной signed family недостаточны}

\subsection{PBW-проекция не сохраняет relation membership покоординатно}

Пусть $F$ --- свободное двухступенно нильпотентное кольцо Ли на
$a<b<c$, а $\Rel$ --- идеал, порождённый $a+c$. Для
\[
 \Theta=\iota(a+c)\iota(b)
\]
PBW-сборка даёт
\[
 \iota(c)\iota(b)=\iota(b)\iota(c)+\iota([c,b]),
 \qquad
 P_1(\Theta)=[c,b].
\]
Но $[c,b]\notin\Rel$. Следовательно,
\[
 \Theta\in\iota(\Rel)U(F)
 \quad\not\Longrightarrow\quad
 P_1(\Theta)\in\Rel.
\]

Аналогично, с двумя spectators можно получить factor-two projection, который
не лежит в образе $\Rel\otimes F\to\Sym^2F$. Поэтому нельзя строить требуемую
цепь по одному relation-left word и надеяться, что primitive correction
автоматически будет целым отношением.

\subsection{Что именно стирается при неправильной абстракции}

Абстрактная family $\mathcal W$ хранит:
\[
 (m_e,\rho_e,X_e)_{e\in E}
 \quad\text{и, следовательно,}\quad
 \Theta_{\mathcal W}.
\]
Она не хранит:

\begin{itemize}
\item какой elementary rewrite породил конкретный descendant;
\item к какому root occurrence он относится;
\item с каким знаком он возник;
\item какой horizontal occurrence парен с каким vertical occurrence;
\item какие relation corrections должны быть сгруппированы до применения
primitive read;
\item какие terminal factor-two rows одновременно задают $\partial z$ и
$\sigma z$.
\end{itemize}

После стирания этих данных агент видел две отдельные задачи:
\[
 \exists z_1:\ \partial z_1=b_E,
 \qquad
 \exists z_2,r:\ \sigma z_2=p_E+r,
\]
но требовалось доказать существование одного $z=z_1=z_2$. Никакая
абстрактная принадлежность идеалу не устанавливает это совпадение.

\begin{dangerbox}
Нельзя сначала выбрать $z$ из factor-two boundary, а затем пытаться отдельно
исправить primitive equation. Нельзя также сначала подобрать primitive lift и
затем доказывать, что его Koszul boundary совпал с $b_E$. Правильный $z$
не выбирается постфактум: он \emph{считывается из конкретных terminal row
occurrences полной трассы}. Поэтому обе координаты совпадают по построению.
\end{dangerbox}

\section{Локальные тождества полной трассы}

\subsection{PBW-transfer обычных факторов}

Для соседних однородных обычных факторов используется буквальное тождество
\begin{equation}\label{eq:ordinary-transfer}
 [\cdots|v|u|\cdots]
 =
 [\cdots|u|v|\cdots]
 +[\cdots|[v,u]|\cdots].
\end{equation}
Первый член уменьшает число PBW-инверсий, второй уменьшает число факторов на
один.

\subsection{Transfer с полным отношением}

Если marked entry является полным отношением $\rho\in\Rel$, то
\begin{equation}\label{eq:relation-transfer}
 [\cdots|v|\rho|\cdots]
 =
 [\cdots|\rho|v|\cdots]
 +[\cdots|[v,\rho]|\cdots].
\end{equation}
Поскольку $\Rel$ --- идеал Ли,
\[
 [v,\rho]\in\Rel.
\]
Следовательно, correction term снова содержит полное отношение. Ни один
однородный хвост не объявляется отношением.

\subsection{Marked transfer и truncation}

Для однородного $v\in F_h$ и $k$-marked relation occurrence
$\langle\rho,k\rangle$, вычисляющегося как $\pr_{\le k}\rho$, имеем
\begin{equation}\label{eq:marked-transfer-doc}
 [\cdots|v|\langle\rho,k\rangle|\cdots]
 =
 [\cdots|\langle\rho,k\rangle|v|\cdots]
 +[\cdots|\langle[v,\rho],k+h\rangle|\cdots].
\end{equation}
Это тождество следует из
\begin{equation}\label{eq:transfer-truncation-doc}
 \pr_{\le k+h}[v,\rho]
 =[v,\pr_{\le k}\rho].
\end{equation}

Шаг truncation имеет вид
\begin{equation}\label{eq:marked-truncation-doc}
 [\cdots|\langle\rho,k\rangle|\cdots]
 =
 [\cdots|\langle\rho,k-1\rangle|\cdots]
 +[\cdots|\rho_k|\cdots].
\end{equation}
Здесь $\rho_k$ --- только $k$-я однородная компонента полного отношения; она
не объявляется элементом $\Rel$. Она является горизонтальным edge quotient
tower. Полный relation label остаётся на первом summand.

\subsection{Критические пары}

Все overlaps локальных правил исчерпываются следующими случаями.

\begin{enumerate}
\item Две непересекающиеся перестановки коммутируют.
\item Три обычных фактора дают обычное тождество Якоби.
\item Overlap с relation entry даёт
\begin{equation}\label{eq:relation-jacobi-doc}
 [u,[v,\rho]]-[v,[u,\rho]]=[[u,v],\rho].
\end{equation}
Все три relation expressions здесь являются полными отношениями.
\item Overlap transfer/truncation замыкается тождеством
\eqref{eq:transfer-truncation-doc}.
\end{enumerate}

Таким образом, локальные клетки образуют действительно замкнутые квадраты:
порядок выполнения соседних шагов меняет только внутренние occurrences, но не
внешнюю signed boundary.

\section{Конечность и правильная мера}

Каждый occurrence обрабатывается детерминированно по лексикографической мере.
Удобно использовать одну из двух эквивалентных конвенций.

\paragraph{Конвенция по active lower weight.}
Если relation occurrence представлен как $\pr_{\ge s}\rho$, используем
\[
 \mu=(p,\operatorname{inv},n+1-s).
\]
Transfer correction уменьшает $p$, обычная перестановка уменьшает
$\operatorname{inv}$, переход к следующему active weight увеличивает $s$ и
уменьшает третью координату.

\paragraph{Конвенция по truncation mark.}
Если occurrence представлен как $\pr_{\le k}\rho$, используем
\[
 \mu'=(p,\operatorname{inv},k).
\]
Marked truncation $k\mapsto k-1$ уменьшает третью координату.

\begin{warning}
В Lean нельзя смешивать эти две записи меры в одном definition. В бумажном
тексте они описывают одну и ту же конечность с противоположно
параметризованной weight coordinate. В формализации следует выбрать только
одну конвенцию; наиболее прямой вариант для правил
\eqref{eq:marked-transfer-doc}--\eqref{eq:marked-truncation-doc} --- мера
$(p,\operatorname{inv},k)$.
\end{warning}

Поскольку $p$, число инверсий и $k$ конечны и неотрицательны, рекурсивный
collector завершается.

\section{Глобальный prefix-ledger invariant}

\subsection{Размеченные occurrences и steps}

Пусть $\mathsf{Occ}$ --- множество \emph{размеченных} row occurrences. Два
occurrences с одинаковым алгебраическим значением, но разным происхождением,
считаются разными базисными элементами. Пусть
\[
 C_1=\Z[\mathsf{Occ}]
\]
--- свободная абелева группа на occurrences.

Пусть $\mathsf{Step}$ --- множество размеченных elementary replacements. Для
шага $s$ заданы один input occurrence $\src(s)$ и конечная signed list outputs
\[
 \out(s)=\sum_j\epsilon_j o_j,
 \qquad \epsilon_j\in\{+1,-1\}.
\]
Положим
\[
 C_2=\Z[\mathsf{Step}],
 \qquad
 \delta(s)=\src(s)-\out(s).
\]

\subsection{Frontier recursion}

Пусть $F_0=E$ --- начальная exceptional signed family с отдельными labels.
Если на шаге $j+1$ обрабатывается occurrence $\src(s_{j+1})$, определим
\begin{equation}\label{eq:frontier-rec}
 F_{j+1}=F_j-\src(s_{j+1})+\out(s_{j+1}).
\end{equation}
Иными словами, input удаляется из frontier, а все signed outputs добавляются.

\begin{lemma}[Prefix ledger]\label{lem:prefix-ledger-doc}
Для каждого $j\ge0$
\begin{equation}\label{eq:prefix-ledger-doc}
 E-F_j=\sum_{i=1}^{j}\delta(s_i).
\end{equation}
\end{lemma}

\begin{proof}
При $j=0$ обе стороны равны нулю. Предположим формулу для $j$. Из
\eqref{eq:frontier-rec} следует
\[
 E-F_{j+1}
 =E-F_j+\src(s_{j+1})-\out(s_{j+1})
 =E-F_j+\delta(s_{j+1}).
\]
Подставляя индукционное предположение, получаем
\[
 E-F_{j+1}
 =\sum_{i=1}^{j+1}\delta(s_i).
\]
\end{proof}

После завершения collector на шаге $N$ положим
\[
 S_E=\sum_{i=1}^{N}s_i\in C_2.
\]
Тогда
\begin{equation}\label{eq:complete-ledger-doc}
 \delta S_E=E-F_N.
\end{equation}
Это точное равенство в свободной абелевой группе на labelled occurrences.
Оно не использует PBW-проекцию, quotient или предварительное сокращение
алгебраически равных terms.

\subsection{Если exceptional family является частью большей трассы}

В полном исходном proof trace начальная family может быть больше $E$. Тогда
нельзя просто удалить прочие корни до построения глобальных диагоналей.
Правильная операция следующая.

\begin{construction}[Descendant-closed exceptional subtrace]
После построения полной размеченной трассы выбираются labelled root
occurrences, составляющие $E$. В подtrace включаются:
\begin{enumerate}
\item все выбранные roots;
\item step, обрабатывающий любой уже включённый occurrence;
\item все outputs такого step;
\item рекурсивно все descendants этих outputs.
\end{enumerate}
\end{construction}

Поскольку labels различают даже равные значения, каждый occurrence имеет
однозначное происхождение. Поэтому descendant closure является конечным
подлесом полной трассы и автоматически удовлетворяет тому же prefix-ledger
identity. Важно, что closure берётся \emph{после} построения полной trace
structure либо одновременно с ней, но никогда не заменяется независимой
wordwise-сборкой каждого значения в $U(F)$.

\section{Additive reads и сохранение локальных тождеств}

Пусть
\[
 \ev:C_1\longrightarrow U(F)
\]
отправляет labelled occurrence в его полное значение. Каждое elementary
replacement является буквальным тождеством, поэтому
\begin{equation}\label{eq:eval-delta-zero}
 \ev\circ\delta=0.
\end{equation}

Нас интересуют два additive reads:
\[
 \beta:C_1\longrightarrow B=\Sym^2A,
 \qquad
 \alpha:C_1\longrightarrow\barF=F/\Rel.
\]
Они задаются как complete factor-two и complete factor-one PBW reads с
последующей нужной проекцией. Из \eqref{eq:eval-delta-zero} следует
\begin{equation}\label{eq:reads-delta-zero}
 \beta\delta=0,
 \qquad
 \alpha\delta=0.
\end{equation}

Слово ``complete'' здесь принципиально: relation leaf сначала собирается как
полный элемент $\Rel$, и лишь затем его класс в $F/\Rel$ объявляется нулём.
Нельзя применять $\alpha$ к отдельным однородным координатам relation tail.

\section{Исчерпывающая классификация terminal frontier}

Рассмотрим однофакторный PBW-вклад веса $n+1$, возникающий при collection из
placed $p$-row. Любой factor-lowering term является bracket tree. По тождеству
Якоби bracket tree целочисленно выражается через сумму left-normed combs.
Поскольку $M=F_{\ge2}$ и $[M,M]=0$, ненулевой метабелев comb содержит не
более одного входа из $M$.

Ниже перечислены все случаи.

\subsection{Единственный $M$-вход есть relation component $\rho_k$}

Все остальные $p-1$ входов должны лежать в $F_1$; иначе появился бы второй
$M$-вход и comb равнялся бы нулю. Вес результата равен
\[
 k+(p-1).
\]
Он равен $n+1$ тогда и только тогда, когда
\begin{equation}\label{eq:diagonal-condition}
 p=n+2-k=m_k.
\end{equation}
Это ровно diagonal cell уровня $k$. Следовательно, ненулевые primitive
contributions такого типа образуют только семейства $H_k$ и $V_k$.

\subsection{Единственный $M$-вход является обычным фактором}

Тогда relation input обязан иметь вес $1$. Пусть он равен $\rho_1$. Если
заменить $\rho_1$ полным отношением
\[
 \rho=\rho_1+\rho_{\ge2},
\]
каждая добавочная компонента $\rho_j$, $j\ge2$, встретится с уже имеющимся
обычным $M$-входом. Полученный comb имеет два входа из $M$ и равен нулю.
Поэтому исходный comb буквально совпадает с comb, содержащим полное
отношение $\rho$, а значит принадлежит $\Rel$ и исчезает после перехода к
$F/\Rel$ или $C$.

\subsection{Все входы имеют вес один}

Для получения веса $n+1$ необходимо $p=n+1$. Снова заменим relation component
$\rho_1$ полным $\rho$. Компонента $\rho_j$, $j\ge2$, после bracket с $n$
входами веса один имеет вес
\[
 j+n\ge n+2,
\]
то есть выше nilpotency class $n+1$, и потому равна нулю. Следовательно,
comb с $\rho_1$ равен полному commutator relation с $\rho$ и также исчезает.

\subsection{Два $M$-входа}

Такой comb равен нулю из-за $[M,M]=0$.

\subsection{Остаётся spectator}

Если после collection остаётся обычный multiplicative spectator, PBW-term
имеет не менее двух факторов. Он не участвует в primitive read $P_1$.

\subsection{Factor number больше $n+1$}

Даже минимальный возможный вес превышает $n+1$, поэтому contribution к
рассматриваемому weight component отсутствует.

\subsection{Factor number два}

Это единственный отдельный terminal case. Он образует oriented quadratic
block и задаёт терминальную цепь $z$.

\begin{proposition}[Полнота frontier classification]\label{prop:frontier-classification}
После применения однофакторного weight-$(n+1)$ read весь terminal frontier
является signed sum следующих и только следующих семейств:
\begin{enumerate}
\item genuine terminal factor-two rows $T$;
\item whole-relation leaves $Q$, исчезающие в $F/\Rel$;
\item horizontal diagonal families $H_k$;
\item vertical diagonal families $V_k$;
\item silent families, read которых равен нулю по одной из причин,
перечисленных выше.
\end{enumerate}
\end{proposition}

\begin{proof}
Каждый primitive contribution после Якоби является суммой combs. Разбиение
по числу $M$-входов, по тому, relation ли единственный $M$-вход, и по числу
оставшихся факторов взаимно исключающе и исчерпывает все combs. В каждом
случае предыдущие подразделы либо относят term к одному из четырёх
несilent семейств, либо доказывают нулевой read. Поэтому иных внешних
семейств нет.
\end{proof}

\section{Диагональное совпадение: локальный расчёт коэффициентов и знаков}

Пусть
\[
 a=\pr_{\le k-1}\rho\in D_{k-1}.
\]
Extension cocycle удовлетворяет
\[
 b_k(a)=-\rho_k\quad\text{в }U_k.
\]
По определению
\[
 \Phi_k(a;x_1\cdots x_{m_k-1})
 =-\Theta_k\bigl(b_k(a);\bar x_1\cdots\bar x_{m_k-1}\bigr),
\]
следовательно,
\begin{equation}\label{eq:horizontal-read-doc}
 \Phi_k(a;x_1\cdots x_{m_k-1})
 =\Theta_k\bigl(\bar\rho_k;\bar x_1\cdots\bar x_{m_k-1}\bigr).
\end{equation}

С другой стороны, в decomposition $A_k=A_{k-1}\oplus F_k$ элемент
$a+\rho_k$ имеет $F_k$-компоненту $\rho_k$. Формула для $T_k$ даёт
\begin{equation}\label{eq:vertical-read-doc}
 T_k\bigl((a+\rho_k)x_1\cdots x_{m_k-1}\bigr)
 =\Theta_k\bigl(\bar\rho_k;\bar x_1\cdots\bar x_{m_k-1}\bigr).
\end{equation}
Сравнивая \eqref{eq:horizontal-read-doc} и
\eqref{eq:vertical-read-doc}, получаем равенство coefficients.

Ориентация boundary клетки даёт horizontal edge со знаком $-$ и vertical
edge со знаком $+$. Поэтому contribution одной labelled diagonal cell равен
\begin{equation}\label{eq:diagonal-contribution-doc}
 -\Phi_k(f_k\chi_k)+T_k(\partial\chi_k).
\end{equation}
Суммирование по labelled occurrences сохраняет исходные multiplicities и
знаки; никакого деления или saturation здесь нет.

\section{Почему терминальные rows определяют один и тот же $z$}

Это центральная точка, которую неправильная абстракция агента уничтожила.

Пусть genuine terminal factor-two frontier имеет вид
\[
 T=\sum_{e\in J}m_e\,t_e,
\]
где каждый $t_e$ содержит relation datum $d_e\in D$ и ordinary datum
$a_e\in A$ с конкретным source placement. Определим
\begin{equation}\label{eq:z-definition}
 z=\sum_{e\in J}m_e\,d_e\otimes a_e\in D\otimes A.
\end{equation}

\subsection{Factor-two coordinate}

Complete factor-two read каждого canonical source-placement row равен
\[
 d_e\odot a_e=\partial(d_e\otimes a_e).
\]
По аддитивности
\begin{equation}\label{eq:beta-T}
 \beta(T)=\partial z.
\end{equation}

\subsection{Primitive coordinate}

По определению source-placement read
\[
 \sigma(d_e\otimes a_e)
 =P_1\bigl(\iota(\lambda d_e)\iota(sa_e)\bigr).
\]
Если в terminal row relation entry первоначально находится не в source
placement, последовательность adjacent relation interchanges переводит её в
source placement. Каждая correction имеет вид
\[
 [v,\lambda d_e]\in\Rel.
\]
Поэтому в $F/\Rel$ primitive read terminal row равен классу
$\sigma(d_e\otimes a_e)$. Суммируя, получаем
\begin{equation}\label{eq:alpha-T}
 \alpha(T)=[\sigma z]\in F/\Rel.
\end{equation}

\begin{successbox}
Равенства \eqref{eq:beta-T} и \eqref{eq:alpha-T} используют один набор
labelled terminal occurrences и одни коэффициенты $m_e$. Поэтому элемент
$z$ в них совпадает по определению. Никакой дополнительной uniqueness lemma
и никакого сравнения двух независимо выбранных lifts не требуется.
\end{successbox}

\section{Построение balanced ledger из полной трассы}

После удаления silent reads, но без удаления их descendants до формирования
глобальной boundary, terminal frontier имеет ориентированное разложение
\begin{equation}\label{eq:frontier-decomp}
 F_N=T-Q-\sum_k(V_k-H_k).
\end{equation}
Знаки здесь выбраны так, чтобы вместе с
$\delta S_E=E-F_N$ получить стандартную balanced-ledger формулу
\begin{equation}\label{eq:balanced-ledger-doc}
 \delta S_E
 =E-T+Q+\sum_k(V_k-H_k).
\end{equation}

Для диагональных семейств локальный closed-square calculation даёт
\begin{equation}\label{eq:reads-diagonal}
 \beta(H_k)=\beta(V_k),
 \qquad
 \alpha(H_k)=\alpha(V_k).
\end{equation}
Первое равенство выражает совместимость factor-two symbol с
transfer/truncation square; второе является detached form chain identity
\[
 \Phi_k(f_k\chi_k)=T_k(\partial\chi_k).
\]
Для whole-relation leaves
\begin{equation}\label{eq:reads-Q}
 \beta(Q)=0,
 \qquad
 \alpha(Q)=0,
\end{equation}
после полного, а не покомпонентного, группирования relation entries.
Наконец, из \eqref{eq:beta-T}--\eqref{eq:alpha-T}
\[
 \beta(T)=\partial z,
 \qquad
 \alpha(T)=[\sigma z].
\]
Таким образом, данные
\[
 (S_E,E,T,Q,(H_k,V_k)_k,z)
\]
являются balanced ledger в точном смысле standalone obstruction formulation.

\begin{theorem}[Balanced ledger из trace]\label{thm:trace-balanced}
Каждая descendant-closed complete marked trace exceptional family $E$
канонически задаёт balanced ledger для $E$.
\end{theorem}

\begin{proof}
Берём $S_E$ как signed sum всех elementary steps descendant-closed subtrace.
Лемма~\ref{lem:prefix-ledger-doc} даёт
$\delta S_E=E-F_N$. Предложение~\ref{prop:frontier-classification}
разлагает $F_N$ в виде \eqref{eq:frontier-decomp}. Терминальная family задаёт
$z$ формулой \eqref{eq:z-definition} и удовлетворяет
\eqref{eq:beta-T}--\eqref{eq:alpha-T}. Whole-relation leaves удовлетворяют
\eqref{eq:reads-Q}, а диагональные families ---
\eqref{eq:reads-diagonal}. Подстановка даёт все поля balanced ledger.
\end{proof}

\section{Из balanced ledger следует нулевой obstruction}

\begin{theorem}[Finite marked Stokes]\label{thm:stokes-doc}
Если $E$ допускает balanced ledger
\eqref{eq:balanced-ledger-doc}, то
\[
 \Obs(E)=0.
\]
\end{theorem}

\begin{proof}
Применим $\beta$ к \eqref{eq:balanced-ledger-doc}. По
\eqref{eq:reads-delta-zero}, $\beta\delta S_E=0$. По
\eqref{eq:reads-Q} relation leaves не дают вклада, а по
\eqref{eq:reads-diagonal} диагональные terms сокращаются. Поэтому
\[
 0=\beta(E)-\beta(T)=b_E-\partial z,
\]
то есть
\begin{equation}\label{eq:first-lift}
 \partial z=b_E.
\end{equation}

Теперь применим $\alpha$. Снова $\alpha\delta S_E=0$, relation leaves
исчезают и диагонали сокращаются. Получаем
\[
 0=\alpha(E)-\alpha(T)=[p_E]-[\sigma z]\quad\text{в }F/\Rel.
\]
Следовательно, существует $r\in\Rel$ такой, что
\begin{equation}\label{eq:second-lift}
 \sigma z=p_E+r.
\end{equation}
Равенства \eqref{eq:first-lift} и \eqref{eq:second-lift} используют один и
тот же $z$, построенный из terminal rows. По критерию
\eqref{eq:obstruction-criterion} имеем $\Obs(E)=0$.
\end{proof}

Теорема~\ref{thm:trace-to-obs} теперь следует из
теорем~\ref{thm:trace-balanced} и \ref{thm:stokes-doc}.

\section{Возврат к terminal discrepancy исходного доказательства}

Полная trace boundary после primitive weight-$(n+1)$ read имеет четыре
несilent family:
\[
 \begin{array}{c|c}
 \text{family} & \text{contribution}\\
 \hline
 \text{terminal one-factor edge}
   & \iota_C(\pr_{1,n+1}\Theta)\\
 \text{horizontal diagonal edge at }k
   & -\Phi_k(f_k\chi_k)\\
 \text{oriented factor-two edge}
   & -Q_n(\chi_n)\\
 \text{vertical diagonal edge at }k
   & +T_k(\partial\chi_k).\\
 \end{array}
\]
Поскольку evaluation каждой elementary step boundary равен нулю, сумма
внешней boundary равна нулю:
\begin{equation}\label{eq:global-external}
 0=
 \iota_C(\pr_{1,n+1}\Theta)
 -\sum_{k=2}^{n-1}\Phi_k(f_k\chi_k)
 -Q_n(\chi_n)
 +\sum_{k=2}^{n-1}T_k(\partial\chi_k).
\end{equation}
По определению \eqref{eq:eps-def} это
\[
 \eps+\sum_{k=2}^{n-1}T_k(\partial\chi_k)=0,
\]
то есть именно \eqref{eq:terminal-discrepancy}.

Остаётся governing vanishing. Функционал $T_k$ поддержан только на
\[
 F_k\Sym^{m_k-1}(F_1)\subseteq\Sym^{m_k}A_k.
\]
Каждый такой monomial имеет total bracket weight
\[
 k+(m_k-1)=k+(n+1-k)=n+1.
\]
Но $\partial\chi_k$ есть complete factor-$m_k$ symbol $\Theta$, а все
многофакторные PBW-коэффициенты $\Theta$ веса не выше $2n$, в частности веса
$n+1$, равны нулю. Следовательно,
\[
 T_k(\partial\chi_k)=0
\]
для всех $2\le k<n$. Из \eqref{eq:terminal-discrepancy} получаем
\[
 \eps=0.
\]
Это завершает проблемную импликацию исходного доказательства.

\section{Особый случай $n=2$}

При $n=2$ промежуточных индексов $2\le k<n$ нет. Target weight равен $3$.
Primitive contribution из трёх factors является целочисленной суммой terms
\[
 [\rho_1,x,y],\qquad x,y\in F_1.
\]
Поскольку $F$ имеет class $3$, для каждой higher component $\rho_j$, $j\ge2$,
\[
 \wt[\rho_j,x,y]=j+2\ge4,
\]
и этот bracket равен нулю. Поэтому
\[
 [\rho_1,x,y]=[\rho,x,y]\in\Rel.
\]
Более трёх positive-weight factors не могут дать вес $3$. Единственный
несilent factor-two contribution есть $Q_2(\chi_2)$. Prefix-ledger identity
даёт
\[
 \iota_C(\pr_{1,3}\Theta)-Q_2(\chi_2)=0.
\]
Таким образом, общий trace argument корректно включает $n=2$, а в Lean этот
случай можно либо оставить в общей конструкции с пустыми sums, либо закрыть
отдельной короткой леммой.

\section{Точная локализация ошибки агента}

\subsection{Неправильная цель}

Агент свёл Step~7 к statement
\[
 \Obs(\mathcal W_{\rm exc})=0
\]
для standalone family и затем пытался доказать его, используя только данные
самой family. Но в исходном аргументе $\mathcal W_{\rm exc}$ не произвольна:
она снабжена происхождением как labelled upper-edge family полной trace.
Правильная формулировка должна быть одной из двух:
\[
 \texttt{CompleteMarkedTrace Wexc -> Obstruction Wexc = 0}
\]
или
\[
 \texttt{trace.exceptionalFamily = Wexc -> Obstruction Wexc = 0}.
\]

\subsection{Потеря lower-context descendants}

При независимой collection relation-left words появляются unwanted primitive
components. В полном trace они компенсируются lower-context diagonal terms.
Если nonexceptional или proper-subset families удалить до образования
глобального signed ledger, компенсация исчезает. Поэтому нельзя передавать в
финальную лемму только очищенный список words без descendant closure.

\subsection{Потеря labels}

Если occurrences отождествлять по алгебраическому значению, нельзя строго
сказать, какой output стал каким последующим input и с каким знаком. Prefix
identity должен жить в свободной группе на labels, а не в $U(F)$. Evaluation
в $U(F)$ применяется только после доказательства combinatorial identity.

\subsection{Попытка построить две цепи вместо одной}

При раздельной работе с factor-two и primitive reads агент неизбежно получает
два потенциальных lifts. Затем вся исходная сложность возникает как
обязанность доказать их совпадение. Полный trace решает это иначе: $z$
определяется ровно один раз из terminal rows, а обе координаты являются
read-функциями одного объекта.

\subsection{Почему разбиение на подпункты переносило сложность}

Если сначала постулировать existence factor-two lift, затем primitive
compatibility, а затем их coherence, то недостающая глобальная trace
перемещается в coherence assumption. Если объединить несколько подпунктов,
агент снова сталкивается с исходной задачей: построить один descendant-closed
ledger. Это не случайность, а признак неверно выбранной границы абстракции.

\section{Какую сессию использовать}

\begin{keybox}
\textbf{Рекомендация: начать новую сессию на чистом коммите.}
Старую сессию следует использовать только для фиксации текущего состояния:
сохранить рабочий diff, записать точное имя незакрытой Lean-теоремы, список
задействованных definitions и последнюю проходящую build-команду. Продолжать
математический поиск в старой сессии не следует: её контекст уже многократно
сжимался и содержит устойчивую неправильную модель задачи --- попытку вывести
obstruction из family без trace witness.
\end{keybox}

Новая сессия должна начинаться не с просьбы ``докажи missing lemma'', а с
конкретной задачи изменить тип данных и провести существующий collector как
proof-producing trace до balanced ledger.

\section{Пошаговый план работы с агентом}

\subsection{Шаг 0. Зафиксировать исходное состояние}

Перед новой сессией:
\begin{enumerate}
\item сохранить текущую рабочую ветку или commit;
\item убедиться, что проект компилируется до единственного известного
placeholder;
\item записать имя terminal theorem и файл, где он находится;
\item сохранить все удачные локальные леммы агента, но удалить или явно
пометить как deprecated ложные abstractions вида ``ideal membership implies
primitive lift'';
\item добавить этот документ в репозиторий рядом с проблемным Lean-файлом.
\end{enumerate}

\subsection{Шаг 1. Провести аудит потери witness}

Первое задание новой сессии: найти, в каком месте Lean-код превращает полную
trace structure в простую signed family. Агент должен ответить кодом, а не
общим рассуждением:
\begin{itemize}
\item какой definition создаёт exceptional family;
\item где теряются labels;
\item где удаляются descendants;
\item где terminal factor-two rows перестают быть доступными для primitive
read;
\item какая существующая функция collection уже может возвращать trace или
может быть минимально расширена.
\end{itemize}
После этого агент сразу вносит минимальное изменение типов; отдельного
долгого отчёта не требуется.

\subsection{Шаг 2. Ввести явный тип complete trace}

Имена ниже схематические; их следует сопоставить существующему коду.

\begin{lstlisting}
structure MarkedOccurrence where
  id        : OccurrenceId
  row       : MarkedRow
  coeff     : Z

structure RewriteStep where
  input     : MarkedOccurrence
  outputs   : List MarkedOccurrence
  validEval : eval input = signedEval outputs
  decreases : measure outputs < measure input

structure CompleteMarkedTrace (E : SignedFamily) where
  steps          : List RewriteStep
  initial        : frontier 0 = E
  next_frontier  : ...
  all_outputs_tracked : ...
  terminal       : isTerminal finalFrontier
\end{lstlisting}

Существенные требования:
\begin{enumerate}
\item occurrence ID входит в равенства frontier;
\item relation field хранит полный $\rho$, mark хранится отдельно;
\item homogeneous component $\rho_k$ является ordinary edge datum, а не
relation proof;
\item outputs каждого step перечислены полностью;
\item collector не возвращает только normal form: он возвращает trace.
\end{enumerate}

\subsection{Шаг 3. Формализовать только три локальных rewrite rules}

Не следует начинать с global obstruction. Сначала должны компилироваться
буквальные леммы:

\begin{lstlisting}
theorem eval_ordinaryTransfer (s : OrdinaryTransfer) :
  eval s.input = signedEval s.outputs := ...

theorem eval_markedRelationTransfer (s : MarkedRelationTransfer) :
  eval s.input = signedEval s.outputs := ...

theorem eval_markedTruncation (s : MarkedTruncation) :
  eval s.input = signedEval s.outputs := ...
\end{lstlisting}

Отдельно:

\begin{lstlisting}
theorem bracket_truncation_square
    (v : Homogeneous F h) (rho : Rel) :
  projLE (k + h) [v, rho] = [v, projLE k rho] := ...
\end{lstlisting}

Ни в одной из этих лемм нельзя предполагать, что $\rho_k\in\Rel$.

\subsection{Шаг 4. Сделать collector proof-producing}

Collector должен быть well-founded recursion по одной фиксированной мере,
например
\[
 (\text{factorCount},\text{unresolvedPairs},\text{mark}).
\]

\begin{lstlisting}
def collectTrace (E : SignedFamily) : CompleteMarkedTrace E :=
  collectAux E (measure E)
termination_by measure E
\end{lstlisting}

Каждый recursive branch обязан вернуть:
\begin{enumerate}
\item конкретный RewriteStep;
\item список descendants;
\item proof local evaluation equality;
\item proof strict decrease;
\item продолжение trace для каждого output, который не terminal.
\end{enumerate}

Если существующий collector уже вычисляет normal form, его не нужно
переписывать с нуля: следует расширить return type журналом steps и доказать,
что забывание журнала даёт прежний result.

\subsection{Шаг 5. Доказать prefix-ledger как общую лемму списка steps}

Эта лемма не должна знать ничего об алгебрах Ли.

\begin{lstlisting}
theorem prefixLedger
    (tr : CompleteMarkedTrace E) (j : Nat) :
  E - tr.frontier j = boundary (tr.steps.take j) := ...

corollary completeLedger (tr : CompleteMarkedTrace E) :
  E - tr.finalFrontier = boundary tr.steps := ...
\end{lstlisting}

Доказательство --- индукция по $j$ с раскрытием одного frontier update. Если
агент здесь начинает использовать PBW, значит типы выбраны неправильно.

\subsection{Шаг 6. Сохранить exceptional family вместе с descendant closure}

Если $E$ является частью более крупной trace, требуется функция:

\begin{lstlisting}
def descendantSubtrace
    (tr : CompleteMarkedTrace I)
    (roots : LabelSubfamily tr.initialFamily) :
    CompleteMarkedTrace roots.family := ...
\end{lstlisting}

Она должна фильтровать по labels и рекурсивной parent relation, а не по
равенству значений rows. Все lower-context descendants остаются внутри.

\subsection{Шаг 7. Классифицировать final frontier}

Вместо попытки сразу доказать \texttt{Obs = 0} определить структуру результата:

\begin{lstlisting}
structure TerminalDecomposition (tr : CompleteMarkedTrace E) where
  terminalRows : SignedFamily
  relationLeaves : SignedFamily
  horizontal : Finset K -> SignedFamily
  vertical   : Finset K -> SignedFamily
  silent     : SignedFamily
  frontier_eq :
    tr.finalFrontier =
      terminalRows - relationLeaves -
      sum k, (vertical k - horizontal k) + silent
  silent_alpha : alpha silent = 0
  silent_beta  : beta silent = 0
\end{lstlisting}

Главная математическая лемма здесь --- comb classification. Её следует
доказывать разбором:
\begin{enumerate}
\item zero, one или at least two $M$-inputs;
\item relation component или ordinary component является единственным
$M$-input;
\item число factors;
\item weight equation $k+p-1=n+1$.
\end{enumerate}

Не следует стирать silent family definitionally: она остаётся в decomposition
и только затем её reads доказываются равными нулю.

\subsection{Шаг 8. Определить общий terminal chain $z$}

\begin{lstlisting}
def terminalChain
    (d : TerminalDecomposition tr) : D tensorProduct Z A :=
  signedSum (d.terminalRows.map rowToTensor)
\end{lstlisting}

Затем две леммы должны быть доказаны для одного definition:

\begin{lstlisting}
theorem beta_terminalRows :
  beta d.terminalRows = partial (terminalChain d) := ...

theorem alpha_terminalRows :
  alpha d.terminalRows = quotientRel (sigma (terminalChain d)) := ...
\end{lstlisting}

Вторая лемма использует только adjacent relation interchanges и факт, что
каждая correction принадлежит полному relation ideal.

\subsection{Шаг 9. Доказать диагональные равенства}

Для каждого $k$:

\begin{lstlisting}
theorem diagonal_beta (k) :
  beta (d.horizontal k) = beta (d.vertical k) := ...

theorem diagonal_alpha (k) :
  alpha (d.horizontal k) = alpha (d.vertical k) := ...
\end{lstlisting}

\texttt{diagonal\_alpha} должна сводиться к уже имеющейся chain transgression
\[
 \Phi_k(f_kc)=T_k(\partial c),
\]
а не заново доказывать глобальную PBW-лемму.

\subsection{Шаг 10. Собрать balanced ledger}

\begin{lstlisting}
def traceToBalancedLedger
    (tr : CompleteMarkedTrace E)
    (d  : TerminalDecomposition tr) :
    BalancedLedger E :=
{ steps := tr.steps
  terminal := d.terminalRows
  relationLeaves := d.relationLeaves
  horizontal := d.horizontal
  vertical := d.vertical
  z := terminalChain d
  ledger_eq := ...
  beta_terminal := beta_terminalRows d
  alpha_terminal := alpha_terminalRows d
  diagonal_beta := diagonal_beta d
  diagonal_alpha := diagonal_alpha d }
\end{lstlisting}

После этого standalone Stokes lemma закрывает obstruction:

\begin{lstlisting}
theorem exceptionalObstructionVanishing
    (tr : CompleteMarkedTrace Wexc) :
    obstruction Wexc = 0 :=
  finiteMarkedStokes (traceToBalancedLedger tr (classifyFrontier tr))
\end{lstlisting}

\subsection{Шаг 11. Подключить к исходной теореме}

Исходный PBW collector должен возвращать одновременно $\chi_k$, $\chi_n$ и
trace witness. Затем terminal discrepancy получается из
\texttt{exceptionalObstructionVanishing} либо непосредственно из evaluated complete
ledger. После этого применяются уже имеющиеся support/governing lemmas,
дающие $T_k(\partial\chi_k)=0$.

\section{Что запрещено агенту делать}

В prompt и review criteria следует явно закрепить следующие запреты.

\begin{enumerate}
\item Не добавлять hypothesis, эквивалентную \texttt{BalancedLedger E} или
\texttt{Obstruction E = 0}.
\item Не использовать \texttt{sorry}, \texttt{admit}, новые \texttt{axiom} или opaque placeholder.
\item Не доказывать wordwise lemma для отдельного relation-left word.
\item Не выводить primitive relation membership только из
$\Theta\in\Rel U(F)$.
\item Не объявлять $\rho_k$ полным отношением.
\item Не удалять nonexceptional, proper-subset или silent descendants до
получения global ledger identity.
\item Не выбирать factor-two lift и primitive lift независимо.
\item Не quotient labels по equality algebraic values.
\item Не заменять полный build локальным toy theorem, который не подключается
к исходному statement.
\item Не менять формулировку исходной математической теоремы и не добавлять к
ней новый public assumption \texttt{trace}; trace должен строиться существующим
collector внутри proof.
\end{enumerate}

\section{Критерии приёмки Lean-решения}

Решение считается завершённым только при одновременном выполнении всех
условий.

\begin{enumerate}
\item Компилируется исходная terminal lemma без placeholders.
\item Компилируется вся исходная theorem, а не только новый standalone файл.
\item В public theorem не добавлены новые assumptions.
\item \texttt{CompleteMarkedTrace} либо существующий эквивалент реально строится из
PBW witness.
\item \texttt{terminalChain} определяется один раз и используется в обеих read
lemmas.
\item Полные отношения сохраняются в relation occurrences; homogeneous tails
не получают relation proof.
\item Prefix-ledger theorem сформулирован на labels.
\item Frontier decomposition содержит явные silent/relation/diagonal
families, а не скрывает их simplifier-ом.
\item Project-wide search не находит новых \texttt{sorry}, \texttt{admit} и \texttt{axiom} в
изменённой части.
\item Проверка аксиом конечной theorem не показывает неожиданных assumptions.
\item Полный project build проходит на чистом checkout текущего commit.
\end{enumerate}

\section{Рекомендуемый prompt для новой сессии Codex}

Ниже дан prompt, который следует передать вместе с этим документом и точными
именами Lean-файлов проекта.

\begin{Verbatim}[fontsize=\small,breaklines=true,breakanywhere=true,frame=single,framesep=2mm]
Работай с текущим компилируемым commit. Главная проблема уже
математически локализована.

НЕПРАВИЛЬНАЯ цель, которую нельзя доказывать:
  SignedRelationWordFamily W -> Obstruction W = 0
или любая её версия, использующая только Theta W in Rel * U(F).
Это утверждение неверно в общем случае.

ПРАВИЛЬНАЯ цель:
  CompleteMarkedTrace Wexc -> BalancedLedger Wexc
  -> Obstruction Wexc = 0.

Wexc не является произвольной family. Она возникает как labelled
upper-edge family полной PBW/truncation trace. Нельзя забывать labels,
parent/descendant data, lower-context descendants или terminal row
placements.

Твоя задача:
1. Найди место, где текущий код стирает complete trace и оставляет
   только Wexc.
2. Минимально расширь collector так, чтобы он возвращал proof-producing
   CompleteMarkedTrace: labelled steps, full outputs, local evaluation
   equalities и termination measure.
3. Докажи generic prefix-ledger identity by induction on steps.
4. Построй descendant-closed subtrace для Wexc, если она является
   частью большей trace.
5. Классифицируй final frontier на terminal rows, whole relation leaves,
   horizontal/vertical diagonals и silent rows.
6. Определи terminalChain z ОДИН РАЗ как signed sum tensors из terminal
   rows.
7. Докажи для этого же z:
      beta terminalRows = partial z
      alpha terminalRows = [sigma z] in F/Rel.
8. Докажи diagonal read equalities, используя существующую
   chain-transgression Phi_k(f_k c) = T_k(partial c).
9. Собери BalancedLedger и примени finiteMarkedStokes.
10. Подключи результат к исходной terminal discrepancy lemma и собери
    весь проект.

Жёсткие ограничения:
- no sorry/admit/new axioms;
- no new public hypotheses in the original theorem;
- rho_k is never treated as a relation;
- do not delete descendants before the global ledger equality;
- do not choose two unrelated chains for the two reads;
- do not replace labelled occurrences by their algebraic values.

После каждого логически завершённого изменения запускай компиляцию
изменённого модуля. Конечный критерий: весь исходный theorem compiles
on a clean checkout and #print axioms contains no unexpected axioms.

Основной математический reference:
terminal_pbw_implication_and_lean_plan_ru.tex,
особенно разделы о prefix ledger, frontier classification,
same terminal chain z и trace-to-balanced-ledger.
\end{Verbatim}

\section{Минимальный и устойчивый варианты патча}

\subsection{Минимальный патч}

Если текущий collector уже создаёт внутренний список шагов, но затем его
забывает, достаточно:
\begin{enumerate}
\item сохранить этот список в return structure;
\item добавить labels;
\item доказать prefix ledger;
\item определить terminal chain по существующему terminal frontier;
\item передать trace witness в missing lemma.
\end{enumerate}
Это предпочтительный путь, если он не требует переписывания основной PBW
инфраструктуры.

\subsection{Устойчивый рефакторинг}

Если collector сразу вычисляет только значения в quotient или normal form,
следует создать отдельный proof-producing collector. Его forgetting map может
воспроизводить старый result. Такой вариант требует больше кода, но устраняет
корневую причину проблемы и делает поздние изменения безопаснее.

\subsection{Правило выбора}

Использовать минимальный патч, если в текущем коде можно указать конкретный
тип, уже содержащий полный step history. Использовать устойчивый рефакторинг,
если history нигде не существует или homogeneous components уже заменили
full relations в типах.

\section{Диагностические исходы}

Новая формализация может обнаружить реальную проблему только в одном из
следующих конкретных мест:

\begin{enumerate}
\item local rewrite identity не типизируется без ложного relation proof для
$\rho_k$;
\item termination measure не убывает на одном из реальных branches;
\item descendant closure теряет output occurrence;
\item frontier classification выдаёт пятый nonsilent primitive case;
\item terminal rows не могут быть поставлены в source placement с correction
в полном $\Rel$;
\item horizontal и vertical reads одной labelled cell имеют разные
coefficients или signs;
\item terminal factor-two и primitive reads требуют разных coefficient lists.
\end{enumerate}

Каждый из этих исходов является точной, проверяемой локализацией. Агент не
должен отвечать общим ``combinatorics is hard''. Он обязан вывести минимальный
непроходящий Lean statement с конкретными terms и types. Однако исходное
бумажное доказательство уже содержит расчёты, исключающие перечисленные
проблемы: local rules буквальны, мера конечна, comb classification исчерпывающа,
диагональные coefficients совпадают, а общий $z$ определяется terminal rows.

\section{Краткий оперативный план для автора}

\begin{enumerate}
\item Зафиксировать текущую ветку и не продолжать старую длинную сессию.
\item Добавить этот .tex и скомпилированный PDF в рабочую ветку.
\item Начать новую сессию с prompt из предыдущего раздела.
\item Потребовать первым кодовым результатом тип \texttt{CompleteMarkedTrace} или
эквивалентное расширение уже существующего collector.
\item Не принимать промежуточную лемму, если она содержит новый assumption,
который по существу равен balanced-ledger existence.
\item Проверить, что \texttt{terminalChain} определён ровно один раз.
\item Проверить descendant closure и labels до анализа PBW reads.
\item После \texttt{traceToBalancedLedger} немедленно подключить Stokes lemma и
исходную terminal theorem.
\item Запустить полный clean build и проверить аксиомы конечной theorem.
\end{enumerate}

\section{Заключение}

Недостающая импликация не требует нового математического принципа. Она
является конечным integral Stokes argument для полной размеченной
PBW/truncation-трассы. Её доказательство состоит из пяти строго отделённых
фактов:

\begin{enumerate}
\item local transfer/truncation steps являются буквальными тождествами и
сохраняют полные отношения;
\item prefix-ledger invariant превращает конечную trace в точное boundary
identity на labels;
\item comb analysis исчерпывающе классифицирует primitive weight-$(n+1)$
frontier;
\item terminal factor-two rows задают один общий $z$, обе координаты которого
получаются двумя reads тех же occurrences;
\item diagonal cells имеют одинаковые coefficients и противоположные boundary
orientations, после чего marked Stokes даёт $\Obs=0$.
\end{enumerate}

Агент застрял потому, что начал после шага, на котором trace witness уже был
забыт. Возвращение этого witness в типы Lean делает правильную импликацию
конструктивной и устраняет необходимость угадывать terminal chain из одного
лишь значения signed family.

\end{document}
